\documentclass[12pt]{ai_conquer2026}
\addbibresource{references.bib}

\title{%
  \begin{center}
        \Large {AI VIET NAM -- AI Conquer 2026}\\[.5ex]  
        \Huge \textbf{Technical Report Template \\ 
        (AI Conquer 2026)}\\[.5ex]
  \end{center}
}

% Author --------------------------------------------------------------------
\posttitle{
    \author{
        \begin{minipage}{\textwidth}
        \centering
        {
        \vspace{-1.5em}
          Tên nhóm / Tên học viên
        }\\
      \end{minipage}
    }
}
\date{}

% Header --------------------------------------------------------------------
\pagestyle{fancy}
\fancyhf{}
\lhead{\href{https://www.facebook.com/aivietnam.edu.vn}{\textcolor{aivnGreen}{\bfseries AI VIET NAM (AI Conquer 2026)}}}
\rhead{\href{https://aivietnam.edu.vn/}{\textcolor{aivnGreen}{\bfseries aivietnam.edu.vn}}}
\fancyfoot[C]{\thepage}
\renewcommand{\headrulewidth}{1.0pt}
\renewcommand{\headrule}{{\color{aivnGreen}\hrule width\headwidth height\headrulewidth \vskip-\headrulewidth}}

\begin{document}
\maketitle
\vspace{-5.5em}
\setlength{\epigraphwidth}{0.6\textwidth}

\section{Tóm tắt báo cáo}

Phần này dùng để tóm tắt ngắn gọn toàn bộ báo cáo kỹ thuật. Nội dung nên cho biết bài toán là gì, dữ liệu được sử dụng như thế nào, hướng giải quyết chính là gì, kết quả nổi bật ra sao, và kết luận quan trọng nhất. Đây là phần giúp người đọc có cái nhìn nhanh trước khi đi vào chi tiết.

\newpage
\setcounter{tocdepth}{2}
\tableofcontents
\thispagestyle{fancy}

\newpage
\section{Giới thiệu}

Phần này giới thiệu ngắn gọn về bài toán, bối cảnh thực hiện và mục tiêu của project. Người viết nên làm rõ vì sao bài toán này đáng quan tâm, project đang hướng tới điều gì, và phạm vi thực hiện nằm ở mức nào. Nếu cần, có thể nêu thêm ứng dụng thực tế hoặc ý nghĩa của bài toán.

\section{Phát biểu bài toán}

Phần này trình bày rõ bài toán mà nhóm đang giải quyết. Người viết nên xác định:
\begin{itemize}
    \item Đầu vào là gì
    \item Đầu ra mong muốn là gì
    \item Loại bài toán là gì (classification, regression, recommendation, generation, v.v.)
    \item Tiêu chí đánh giá chính là gì
\end{itemize}

Phần này càng rõ thì các phần sau càng dễ theo dõi.

\section{Dữ liệu}

Phần này mô tả dữ liệu được sử dụng trong project. Nội dung có thể bao gồm:
\begin{itemize}
    \item Nguồn dữ liệu
    \item Kích thước dữ liệu
    \item Các trường thông tin chính
    \item Những vấn đề ban đầu của dữ liệu
    \item Các nhận xét quan trọng sau khi khảo sát dữ liệu
\end{itemize}

Nếu có hình minh họa, bảng thống kê hoặc ví dụ mẫu dữ liệu thì có thể chèn thêm ở phần này.

\section{Tiền xử lý dữ liệu}

Phần này trình bày các bước xử lý dữ liệu trước khi đưa vào mô hình hoặc hệ thống. Nội dung có thể bao gồm:
\begin{itemize}
    \item Xử lý missing values
    \item Loại bỏ dữ liệu trùng hoặc lỗi
    \item Chuẩn hóa định dạng
    \item Encoding / tokenization / normalization
    \item Feature engineering
    \item Chia train / validation / test
\end{itemize}

Người viết nên giải thích vì sao cần các bước đó và đầu ra sau tiền xử lý là gì.

\section{Phương pháp / Cách tiếp cận}

Phần này mô tả hướng giải quyết của nhóm. Người viết nên trình bày:
\begin{itemize}
    \item Ý tưởng tổng quát
    \item Lý do chọn cách tiếp cận này
    \item Baseline được sử dụng là gì
    \item Những phương pháp hoặc mô hình được thử nghiệm
\end{itemize}

Nếu project không dùng model học máy mà đi theo hướng rule-based, logic system hoặc pipeline xử lý thì phần này vẫn cần làm rõ cách tiếp cận tổng thể.

\section{Triển khai hệ thống}

Phần này mô tả cách nhóm hiện thực hóa giải pháp. Nội dung có thể bao gồm:
\begin{itemize}
    \item Cấu trúc hệ thống
    \item Các thành phần chính trong pipeline
    \item Công cụ và thư viện sử dụng
    \item Cách các thành phần kết nối với nhau
    \item Các quyết định kỹ thuật quan trọng
\end{itemize}

Nếu phù hợp, có thể bổ sung sơ đồ pipeline hoặc mô tả luồng xử lý từ input đến output.

\section{Thực nghiệm và đánh giá}

Phần này trình bày cách nhóm kiểm tra và đánh giá giải pháp. Nội dung có thể bao gồm:
\begin{itemize}
    \item Thiết lập thực nghiệm
    \item Metric sử dụng
    \item Các mô hình hoặc phương án được so sánh
    \item Kết quả đạt được
    \item Nhận xét về kết quả
\end{itemize}

Nếu có bảng kết quả hoặc biểu đồ thì nên trình bày rõ ràng ở phần này.

\section{Phân tích lỗi và thảo luận}

Phần này giúp báo cáo có chiều sâu hơn. Người viết nên trình bày:
\begin{itemize}
    \item Các trường hợp hệ thống làm tốt
    \item Các trường hợp hệ thống còn lỗi
    \item Nguyên nhân có thể dẫn đến lỗi
    \item Hạn chế hiện tại của cách làm
\end{itemize}

Đây là phần rất quan trọng vì nó cho thấy nhóm không chỉ chạy ra kết quả, mà còn hiểu được kết quả đó.

\section{Kết luận và hướng phát triển}

Phần cuối dùng để tổng kết lại toàn bộ project. Nội dung nên trả lời ngắn gọn các câu hỏi:
\begin{itemize}
    \item Nhóm đã giải quyết bài toán bằng cách nào
    \item Kết quả chính đạt được là gì
    \item Hạn chế lớn nhất hiện tại là gì
    \item Nếu có thêm thời gian thì nên cải thiện theo hướng nào
\end{itemize}

\newpage
\section{Tài liệu tham khảo}
\label{section:references}
\nocite{*}
\vspace{-3.2em}
\printbibliography

\newpage
\begin{appendices}

\section{Phụ lục}

Phần này có thể dùng để bổ sung các nội dung như:
\begin{itemize}
    \item Ví dụ dữ liệu
    \item Bảng kết quả đầy đủ
    \item Mô tả thêm về tham số
    \item Hình ảnh minh họa
    \item Link GitHub / Demo / Video
\end{itemize}

\end{appendices}

\end{document}